% Part: first-order-logic
% Chapter: models-theories
% Section: set-theory

\documentclass[../../../include/open-logic-section]{subfiles}

\begin{document}

\olfileid{fol}{mat}{set}

\olsection{نظرية المجموعات}

تكاد نظرية المجموعات تتسع لبناء الرياضيات كلها. ولهذا البناء
مطالب، أولها وضع بديهيات للعلاقة~$\in$. وقد وُضعت لها أنساق
متعددة، تتعارض أحيانًا في بعض خواص~$\in$. وأشهرها
$\Log{ZFC}$، أي نظرية زرميلو--فرينكل مع بديهية الاختيار؛
وهي أوسع تلك الأنساق استعمالًا ودراسة، إذ تكفي بديهياتها لإثبات
جلّ ما يتوقع الرياضيون إثباته. غير أن البحث في ذلك مسبوق بسؤال:
كيف \emph{نعبّر} في هذه اللغة عما نريده من مسائل الرياضيات؟
فليس في اللغة !!{constant}s ولا !!{function}s، فلا يظهر أولًا
كيف نسمّي مجموعات بعينها، مثل $\emptyset$ و$\Nat$، أو نصف
عمليات المجموعات، مثل ${\OLId{latin-looped}{X}} \cup {\OLId{latin-looped}{Y}}$ و$\Pow{{\OLId{latin-looped}{X}}}$، فضلًا عن
التراكيب الأخرى، كالعلاقات والدوال، التي تبدو مشتملة على غير
المجموعات.

ولنبدأ بعلاقة «هو مجموعة جزئية من»، فإنها تقارب علاقة «هو عنصر
من» في الأهمية. نستطيع \emph{تعريف} الأولى بالثانية، وذلك بأن
نجد !!a{formula}~$!A({\OLId{latin-ordinary}{x}}, {\OLId{latin-ordinary}{y}})$ من لغة نظرية المجموعات يرضيها
الزوج~$\tuple{{\OLId{latin-looped}{X}}, {\OLId{latin-looped}{Y}}}$ إذا وفقط إذا كان ${\OLId{latin-looped}{X}} \subseteq {\OLId{latin-looped}{Y}}$.
وإنما تكون ${\OLId{latin-looped}{X}}$ مجموعة جزئية من ${\OLId{latin-looped}{Y}}$ متى كانت كل
!!{element}s~${\OLId{latin-looped}{X}}$ !!{element}s في~${\OLId{latin-looped}{Y}}$ أيضًا. فتعريف
$\subseteq$ هو الصيغة
\[
\lforall[{\OLId{latin-ordinary}{z}}][({\OLId{latin-ordinary}{z}} \in {\OLId{latin-ordinary}{x}} \lif {\OLId{latin-ordinary}{z}} \in {\OLId{latin-ordinary}{y}})]
\]
فحيث نريد العلاقة~$\subseteq$ في صيغة، نضع مكانها هذه العبارة،
مع الإحلال المناسب محل ${\OLId{latin-ordinary}{x}}$ و${\OLId{latin-ordinary}{y}}$، وتغيير اسم المتغير المقيد~${\OLId{latin-ordinary}{z}}$
إذا احتيج إليه. ومن أمثلة ذلك امتدادية المجموعات: إذا كانت
مجموعتان~${\OLId{latin-ordinary}{x}}$ و${\OLId{latin-ordinary}{y}}$، كل واحدة منهما محتواة في الأخرى، فهما مجموعة
واحدة. نكتب ذلك
$\lforall[{\OLId{latin-ordinary}{x}}][\lforall[{\OLId{latin-ordinary}{y}}][(({\OLId{latin-ordinary}{x}} \subseteq {\OLId{latin-ordinary}{y}} \land {\OLId{latin-ordinary}{y}}
    \subseteq {\OLId{latin-ordinary}{x}}) \lif {\OLId{latin-ordinary}{x}} = {\OLId{latin-ordinary}{y}})]]$؛ فإذا رددنا $\subseteq$
إلى التعريف المتقدم حصل
\[
\lforall[{\OLId{latin-ordinary}{x}}][\lforall[{\OLId{latin-ordinary}{y}}][((\lforall[{\OLId{latin-ordinary}{z}}][({\OLId{latin-ordinary}{z}} \in {\OLId{latin-ordinary}{x}} \lif {\OLId{latin-ordinary}{z}} \in {\OLId{latin-ordinary}{y}})] \land
    \lforall[{\OLId{latin-ordinary}{z}}][({\OLId{latin-ordinary}{z}} \in {\OLId{latin-ordinary}{y}} \lif {\OLId{latin-ordinary}{z}} \in {\OLId{latin-ordinary}{x}})]) \lif {\OLId{latin-ordinary}{x}} = {\OLId{latin-ordinary}{y}})]].
\]
وهذه هي «بديهية الامتدادية»، إحدى بديهيات~$\Log{ZFC}$.

ليس في اللغة !!{constant} لـ~$\emptyset$؛ غير أن عبارة «${\OLId{latin-ordinary}{x}}$
خالية» تؤدى بـ~$\lnot \lexists[{\OLId{latin-ordinary}{y}}][{\OLId{latin-ordinary}{y}} \in {\OLId{latin-ordinary}{x}}]$. فتكون عبارة
«$\emptyset$ موجودة» ال!!{sentence}~$\lexists[{\OLId{latin-ordinary}{x}}][\lnot
\lexists[{\OLId{latin-ordinary}{y}}][{\OLId{latin-ordinary}{y}} \in {\OLId{latin-ordinary}{x}}]]$. وهذه بديهية أخرى من بديهيات~$\Log{ZFC}$.
(لاحظ أن بديهية الامتدادية تقتضي وجود مجموعة خالية واحدة فقط.) ومتى
أردنا ذكر $\emptyset$ في لغة نظرية المجموعات، كتبنا ذلك على
صورة «توجد مجموعة خالية و\dots». وللتعبير، مثلًا، عن أن $\emptyset$
مجموعة جزئية من كل مجموعة، يمكننا أن نكتب
\[
\lexists[{\OLId{latin-ordinary}{x}}][(\lnot\lexists[{\OLId{latin-ordinary}{y}}][{\OLId{latin-ordinary}{y}} \in {\OLId{latin-ordinary}{x}}] \land \lforall[{\OLId{latin-ordinary}{z}}][{\OLId{latin-ordinary}{x}} \subseteq
    {\OLId{latin-ordinary}{z}}])]
\]
على أن ${\OLId{latin-ordinary}{x}} \subseteq {\OLId{latin-ordinary}{z}}$ يجب، بالطبع، أن يستبدل بدوره بتعريفه.

وفي عمليات المجموعات، مثل ${\OLId{latin-looped}{X}} \cup {\OLId{latin-looped}{Y}}$ و$\Pow{{\OLId{latin-looped}{X}}}$، نسلك
مسلكًا شبيهًا بذلك. فلا رموز دوال في لغة نظرية المجموعات، لكننا
نستطيع التعبير عن العلاقتين الداليتين ${\OLId{latin-looped}{X}} \cup {\OLId{latin-looped}{Y}} = {\OLId{latin-looped}{Z}}$ و$\Pow{{\OLId{latin-looped}{X}}} = {\OLId{latin-looped}{Y}}$
بالصيغتين
\begin{align*}
& \lforall[{\OLId{latin-ordinary}{u}}][(({\OLId{latin-ordinary}{u}} \in {\OLId{latin-ordinary}{x}} \lor {\OLId{latin-ordinary}{u}} \in {\OLId{latin-ordinary}{y}}) \liff {\OLId{latin-ordinary}{u}} \in {\OLId{latin-ordinary}{z}})]\\
& \lforall[{\OLId{latin-ordinary}{u}}][({\OLId{latin-ordinary}{u}} \subseteq {\OLId{latin-ordinary}{x}} \liff {\OLId{latin-ordinary}{u}} \in {\OLId{latin-ordinary}{y}} )]
\end{align*}
لأن !!{element}s ${\OLId{latin-looped}{X}} \cup {\OLId{latin-looped}{Y}}$ هي بالضبط المجموعات التي تكون
!!{element}s في~${\OLId{latin-looped}{X}}$ أو !!{element}s في~${\OLId{latin-looped}{Y}}$، و!!{element}s
$\Pow{{\OLId{latin-looped}{X}}}$ هي بالضبط المجموعات الجزئية من~${\OLId{latin-looped}{X}}$. ولا يلزم من ذلك جواز
استعمال ${\OLId{latin-ordinary}{x}} \cup {\OLId{latin-ordinary}{y}}$ أو $\Pow{{\OLId{latin-ordinary}{x}}}$ كما لو كانا حدين؛ إذ لا يمكننا إلا
استعمال ال!!{formula}s الكاملة التي تعرّف العلاقتين ${\OLId{latin-looped}{X}} \cup {\OLId{latin-looped}{Y}} = {\OLId{latin-looped}{Z}}$
و$\Pow{{\OLId{latin-looped}{X}}} = {\OLId{latin-looped}{Y}}$. ومجرد التعريف لا يثبت أن هاتين العلاقتين تُرضيان، ولا أن الاتحادات
ومجموعات القوى موجودة لكل ما نختاره من المجموعات. فمثلًا،
ال!!{sentence} $\lforall[{\OLId{latin-ordinary}{x}}][\lexists[{\OLId{latin-ordinary}{y}}][\Pow{{\OLId{latin-ordinary}{x}}} = {\OLId{latin-ordinary}{y}}]]$ بديهية أخرى
من بديهيات~$\Log{ZFC}$ (بديهية مجموعة القوى).

وأما الأزواج المرتبة والدوال، فلا بد من بيان كيف تُجعل أنواعًا
من المجموعات. فمن تعريفات الزوج المرتب $\tuple{{\OLId{latin-ordinary}{x}}, {\OLId{latin-ordinary}{y}}}$ أن يكون
المجموعة $\{\{{\OLId{latin-ordinary}{x}}\}, \{{\OLId{latin-ordinary}{x}}, {\OLId{latin-ordinary}{y}}\}\}$. وليس لنا في أصل اللغة، كما
تقدم، أن نضيف !!a{function} تسمّي هذه المجموعة، بل نعرّف العلاقة
$\tuple{{\OLId{latin-ordinary}{x}}, {\OLId{latin-ordinary}{y}}} = {\OLId{latin-ordinary}{z}}$، أي $\{\{{\OLId{latin-ordinary}{x}}\}, \{{\OLId{latin-ordinary}{x}}, {\OLId{latin-ordinary}{y}}\}\} = {\OLId{latin-ordinary}{z}}$، بالصيغة
\[
\lforall[{\OLId{latin-ordinary}{u}}][({\OLId{latin-ordinary}{u}} \in {\OLId{latin-ordinary}{z}} \liff (\lforall[{\OLId{latin-ordinary}{v}}][({\OLId{latin-ordinary}{v}} \in {\OLId{latin-ordinary}{u}} \liff {\OLId{latin-ordinary}{v}} = {\OLId{latin-ordinary}{x}})] \lor
  \lforall[{\OLId{latin-ordinary}{v}}][({\OLId{latin-ordinary}{v}} \in {\OLId{latin-ordinary}{u}} \liff ({\OLId{latin-ordinary}{v}} = {\OLId{latin-ordinary}{x}} \lor {\OLId{latin-ordinary}{v}} = {\OLId{latin-ordinary}{y}}))]))]
\]
وتقول هذه الصيغة إن !!{element}s~${\OLId{latin-ordinary}{u}}$ في~${\OLId{latin-ordinary}{z}}$ هي بالضبط المجموعات التي
يكون !!{element}ها الوحيد هو ${\OLId{latin-ordinary}{x}}$، أو تكون !!{element}sها الوحيدة هي
${\OLId{latin-ordinary}{x}}$ و~${\OLId{latin-ordinary}{y}}$ (وبعبارة
أخرى، المجموعات المطابقة إما لـ~$\{{\OLId{latin-ordinary}{x}}\}$ وإما لـ~$\{{\OLId{latin-ordinary}{x}}, {\OLId{latin-ordinary}{y}}\}$). وبه
يمكن التعبير عن علاقات أخرى، منها ${\OLId{latin-looped}{X}} \times {\OLId{latin-looped}{Y}} = {\OLId{latin-looped}{Z}}$:
\[
\lforall[{\OLId{latin-ordinary}{z}}][({\OLId{latin-ordinary}{z}} \in {\OLId{latin-looped}{Z}} \liff \lexists[{\OLId{latin-ordinary}{x}}][\lexists[{\OLId{latin-ordinary}{y}}][({\OLId{latin-ordinary}{x}} \in
        {\OLId{latin-looped}{X}} \land {\OLId{latin-ordinary}{y}} \in {\OLId{latin-looped}{Y}} \land \tuple{{\OLId{latin-ordinary}{x}}, {\OLId{latin-ordinary}{y}}} = {\OLId{latin-ordinary}{z}})]])]
\]

وللدالة ${\OLId{latin-ordinary}{f}} \colon {\OLId{latin-looped}{X}} \to {\OLId{latin-looped}{Y}}$ أن تُمثّل بالعلاقة ${\OLId{latin-ordinary}{f}}({\OLId{latin-ordinary}{x}})
= {\OLId{latin-ordinary}{y}}$، أي بمجموعة الأزواج~$\Setabs{\tuple{{\OLId{latin-ordinary}{x}},{\OLId{latin-ordinary}{y}}}}{{\OLId{latin-ordinary}{f}}({\OLId{latin-ordinary}{x}}) = {\OLId{latin-ordinary}{y}}}$.
فتكون المجموعة~${\OLId{latin-ordinary}{f}}$ دالة من ${\OLId{latin-looped}{X}}$ إلى ${\OLId{latin-looped}{Y}}$ بثلاثة شروط: أولها أن
تكون علاقة $\subseteq {\OLId{latin-looped}{X}} \times {\OLId{latin-looped}{Y}}$؛ وثانيها الكلية، أي أن لكل
${\OLId{latin-ordinary}{x}} \in {\OLId{latin-looped}{X}}$ بعض ${\OLId{latin-ordinary}{y}} \in {\OLId{latin-looped}{Y}}$ يحقق $\tuple{{\OLId{latin-ordinary}{x}}, {\OLId{latin-ordinary}{y}}} \in {\OLId{latin-ordinary}{f}}$؛
وثالثها الدالية، أي إن $\tuple{{\OLId{latin-ordinary}{x}}, {\OLId{latin-ordinary}{y}}}, \tuple{{\OLId{latin-ordinary}{x}}, {\OLId{latin-ordinary}{y}}'} \in {\OLId{latin-ordinary}{f}}$
يقتضي ${\OLId{latin-ordinary}{y}} = {\OLId{latin-ordinary}{y}}'$، فلا تتعدد القيمة. ولهذا تُكتب عبارة «${\OLId{latin-ordinary}{f}}$ دالة
من ${\OLId{latin-looped}{X}}$ إلى ${\OLId{latin-looped}{Y}}$» هكذا:
\begin{align*}
 \lforall[{\OLId{latin-ordinary}{u}}][({\OLId{latin-ordinary}{u}} \in {\OLId{latin-ordinary}{f}} \lif {}] & \lexists[{\OLId{latin-ordinary}{x}}][\lexists[{\OLId{latin-ordinary}{y}}][({\OLId{latin-ordinary}{x}} \in {\OLId{latin-looped}{X}} \land {\OLId{latin-ordinary}{y}} \in
      {\OLId{latin-looped}{Y}} \land \tuple{{\OLId{latin-ordinary}{x}}, {\OLId{latin-ordinary}{y}}} = {\OLId{latin-ordinary}{u}})]]) \land {}\\
 \lforall[{\OLId{latin-ordinary}{x}}][({\OLId{latin-ordinary}{x}} \in {\OLId{latin-looped}{X}} \lif {}] &
      (\lexists[{\OLId{latin-ordinary}{y}}][({\OLId{latin-ordinary}{y}} \in {\OLId{latin-looped}{Y}} \land \mathrm{maps}({\OLId{latin-ordinary}{f}}, {\OLId{latin-ordinary}{x}}, {\OLId{latin-ordinary}{y}}))] \land {}\\
& (\lforall[{\OLId{latin-ordinary}{y}}][\lforall[{\OLId{latin-ordinary}{y}}'][((\mathrm{maps}({\OLId{latin-ordinary}{f}}, {\OLId{latin-ordinary}{x}}, {\OLId{latin-ordinary}{y}}) \land
    \mathrm{maps}({\OLId{latin-ordinary}{f}}, {\OLId{latin-ordinary}{x}}, {\OLId{latin-ordinary}{y}}')) \lif {\OLId{latin-ordinary}{y}} = {\OLId{latin-ordinary}{y}}')]]))
\end{align*}
حيث تختصر $\mathrm{maps}({\OLId{latin-ordinary}{f}}, {\OLId{latin-ordinary}{x}}, {\OLId{latin-ordinary}{y}})$ الصيغة $\lexists[{\OLId{latin-ordinary}{v}}][({\OLId{latin-ordinary}{v}} \in {\OLId{latin-ordinary}{f}} \land
  \tuple{{\OLId{latin-ordinary}{x}}, {\OLId{latin-ordinary}{y}}} = {\OLId{latin-ordinary}{v}})]$ (وهذه ال!!{formula} تعبر عن «${\OLId{latin-ordinary}{f}}({\OLId{latin-ordinary}{x}}) = {\OLId{latin-ordinary}{y}}$»).

وبعد هذا يسهل التعبير، مثلًا، عن كون ${\OLId{latin-ordinary}{f}}\colon {\OLId{latin-looped}{X}} \to {\OLId{latin-looped}{Y}}$
!!{injective}:
\begin{multline*}
 {\OLId{latin-ordinary}{f}} \colon {\OLId{latin-looped}{X}} \to {\OLId{latin-looped}{Y}} \land \lforall[{\OLId{latin-ordinary}{x}}][\lforall[{\OLId{latin-ordinary}{x}}'][(({\OLId{latin-ordinary}{x}} \in {\OLId{latin-looped}{X}} \land {\OLId{latin-ordinary}{x}}' \in
     {\OLId{latin-looped}{X}} \land {}]] \\
 \lexists[{\OLId{latin-ordinary}{y}}][(\mathrm{maps}({\OLId{latin-ordinary}{f}}, {\OLId{latin-ordinary}{x}}, {\OLId{latin-ordinary}{y}}) \land \mathrm{maps}({\OLId{latin-ordinary}{f}},
        {\OLId{latin-ordinary}{x}}', {\OLId{latin-ordinary}{y}}))]) \lif {\OLId{latin-ordinary}{x}} = {\OLId{latin-ordinary}{x}}')
\end{multline*}
تكون الدالة~${\OLId{latin-ordinary}{f}}\colon {\OLId{latin-looped}{X}} \to {\OLId{latin-looped}{Y}}$ !!{injective} إذا وفقط إذا كان
${\OLId{latin-ordinary}{x}} = {\OLId{latin-ordinary}{x}}'$ كلما أرسلت ${\OLId{latin-ordinary}{x}}, {\OLId{latin-ordinary}{x}}' \in {\OLId{latin-looped}{X}}$ إلى قيمة واحدة~${\OLId{latin-ordinary}{y}}$. وإذا اختصرنا
هذه الصيغة إلى $\mathrm{inj}({\OLId{latin-ordinary}{f}}, {\OLId{latin-looped}{X}}, {\OLId{latin-looped}{Y}})$، أمكننا أن نصوغ بلغة نظرية المجموعات نتيجة غير يسيرة، هي مبرهنة كانتور: لا توجد
دالة !!{injective} من $\Pow{{\OLId{latin-looped}{X}}}$ إلى ${\OLId{latin-looped}{X}}$:
\[
\lforall[{\OLId{latin-looped}{X}}][\lforall[{\OLId{latin-looped}{Y}}][(\Pow{{\OLId{latin-looped}{X}}} = {\OLId{latin-looped}{Y}} \lif
    \lnot\lexists[{\OLId{latin-ordinary}{f}}][\mathrm{inj}({\OLId{latin-ordinary}{f}}, {\OLId{latin-looped}{Y}}, {\OLId{latin-looped}{X}})])]]
\]

وقد يبدو أن نظرية المجموعات تحتاج إلى بديهية تجعل لكل خاصية
معرِّفة مجموعة. فإذا كانت $!A({\OLId{latin-ordinary}{x}})$ صيغة في لغة المجموعات
فيها المتغير~${\OLId{latin-ordinary}{x}}$ حرًا، أمكن النظر في ال!!{sentence}
\[
\lexists[{\OLId{latin-ordinary}{y}}][\lforall[{\OLId{latin-ordinary}{x}}][({\OLId{latin-ordinary}{x}} \in {\OLId{latin-ordinary}{y}} \liff !A({\OLId{latin-ordinary}{x}}))]].
\]
وتقرر هذه ال!!{sentence} أن ثمة مجموعة~${\OLId{latin-ordinary}{y}}$ تكون !!{element}sها كل
قيم ${\OLId{latin-ordinary}{x}}$ التي ترضي~$!A({\OLId{latin-ordinary}{x}})$ وتلك القيم وحدها. ويسمى هذا المخطط «مبدأ
الاستيعاب». ومع ما يبدو فيه من النفع، فهو غير متسق. خذ
$!A({\OLId{latin-ordinary}{x}}) \ident \lnot {\OLId{latin-ordinary}{x}} \in {\OLId{latin-ordinary}{x}}$؛ عندئذ يقرر مبدأ الاستيعاب
\[
\lexists[{\OLId{latin-ordinary}{y}}][\lforall[{\OLId{latin-ordinary}{x}}][({\OLId{latin-ordinary}{x}} \in {\OLId{latin-ordinary}{y}} \liff {\OLId{latin-ordinary}{x}} \notin {\OLId{latin-ordinary}{x}})]],
\]
أي إنه يقرر وجود مجموعة جميع المجموعات غير المنتمية إلى نفسها. ويمتنع
وجود مثل هذه المجموعة؛ وتلك مفارقة راسل. والواقع أن
$\Log{ZFC}$ تحتوي نسخة مقيدة---ومتسقة---من هذا المبدأ، هي مبدأ الفصل:
\[
\lforall[{\OLId{latin-ordinary}{z}}][\lexists[{\OLId{latin-ordinary}{y}}][\lforall[{\OLId{latin-ordinary}{x}}][({\OLId{latin-ordinary}{x}} \in {\OLId{latin-ordinary}{y}} \liff ({\OLId{latin-ordinary}{x}} \in {\OLId{latin-ordinary}{z}} \land
      !A({\OLId{latin-ordinary}{x}})))]]].
\]

\begin{prob}
بين أن مبدأ الاستيعاب غير متسق بإعطاء !!a{derivation} يبين أن
\[
\lexists[{\OLId{latin-ordinary}{y}}][\lforall[{\OLId{latin-ordinary}{x}}][({\OLId{latin-ordinary}{x}} \in {\OLId{latin-ordinary}{y}} \liff {\OLId{latin-ordinary}{x}} \notin {\OLId{latin-ordinary}{x}})]] \Proves \lfalse.
\]
وللاستعانة، أثبت أولًا أن $({\OLId{latin-looped}{A}} \lif \lnot {\OLId{latin-looped}{A}}) \land (\lnot {\OLId{latin-looped}{A}} \lif {\OLId{latin-looped}{A}})
\Proves \lfalse$.
\end{prob}
\end{document}
